\chapter*{Remerciements}

Bien que ce doctorat ait commencé il y a quatre ans, ses fondations ont été érigées, je crois, au cours de ma vie entière. Sans l'apport de chaque rencontre, chaque discussion, chaque désaccord, chaque encouragement, chaque sanction, il est probable, pour ne pas dire certain, que cette aventure ne soit pas parvenue à bon port. Je me dois donc de remercier mes collègues, mes amis et ma famille, bien sûr, mais aussi tous ceux et celles qui, à la manière des clous d'une planche de Galton, ont dévié la trajectoire de ma vie.

Avant tout, je souhaite remercier mon superviseur Gilles Louppe, sans qui mon doctorat aurait été littéralement impossible. Blague à part, Gilles, c'est à toi que je dois ma passion pour l'IA, mon goût pour la recherche, et tout ce que je sais de la rigueur et de l'ambition scientifiques. J'ai dit par le passé que tu étais avare de compliments. Je ne le comprends que maintenant, mais la liberté et la confiance que tu m'as accordées étaient autant d'approbations silencieuses. Pour ça et tout le reste, merci.

Je suis reconnaissant aux membres du jury, Marc Bocquet, Pierre Geurts, Philipp Hennig, Antoine Wehenkel et Louis Wehenkel, pour l'intérêt qu'ils ont porté à mon travail et pour avoir pris le temps de lire cette dissertation. J'ai grandement apprécié leurs commentaires constructifs, ainsi que leur engagement pendant la défense. Je souhaite aussi remercier Pierre pour son intégrité en tant qu'enseignant et Antoine pour ses précieux conseils tout au long de mon doctorat.

Presque dix années se sont déroulées entre mon entrée à l'Université de Liège et ma sortie. Mes collègues de Montefiore ont rendu les quatre dernières particulièrement mémorables. Merci Gérôme et Omer, d'avoir partagé mon bureau, mes idées et mes bêtises. Merci Arnaud, Sacha, Victor, Adrien, Lénea, Élise, Fanny, Thomas, Sacha et Yakup, d'avoir contribué à un Grouppe dynamique et stimulant. Merci Pascal, Matthias, Yann, Gaspard, Sammy, Victor, Adrien, Jérôme et Renaud, pour la bonne humeur et les distractions, désirées ou imposées. Merci Anna, Florent, Sven, Astrid, Cyril, Laura, Robin, Simon et Vicole, pour les pauses de midi débridées entre habitants de Thiercelieux.

I was also lucky enough to leave Montefiore for a four-month internship in the heart of Manhattan. Thank you, Shirley and François, for this fantastic opportunity which I will never forget. Thank you, Ruben, Mike, and Rudy, for welcoming me into the team and guiding me in my work. Finally, thank you to my roommates, Lucas and Samuele, for sharing those thirty square meters, questionable jokes, and endless debates. I believe Gauss, who said \enquote{Mathematics is the queen of the sciences}, would have enjoyed playing Minecraft.

Si j'ai pu persévérer dans cette aventure, c'est en partie grâce à mes amis, qui m'ont permis de m'évader de cette addiction qu'est la recherche. Merci Nikolay et Pedro, de rester mes plus proches amis, malgré la distance. Merci Maxime, d'avoir mêlé code et délires avec moi pendant presque quinze ans. Merci Simon, Louis et Ivan, d'être mes partenaires de jeu, que ce soit autour d'une table ou derrière un clavier. Merci Adrien, Bastien, Justin, Quentin, Thomas, Samuel et Valentin, pour ces soirées (in)oubliables et ces matchs légendaires.

Finalement, ma famille est sans doute l'ingrédient principal de ma réussite. Merci Papi, de m'avoir mis des outils entre les mains et des étincelles dans les yeux dès mon plus jeune âge. Merci Nonna, de m'avoir toujours accueilli et nourri à bras ouverts. Merci Lino, Juan, Sophie, Quentin et Caroline, de me remettre les pieds sur terre quand il le faut. Merci Léo, d'être mon anti-dépresseur naturel. Merci Mathilde, d'avoir été ma bouée de sauvetage tout au long de ma scolarité. Merci David, d'être bien plus qu'un beau-frère rigolo. Merci Marion, d'être toujours présente et la meilleure des conseillères. Merci Maman, de penser à tout avant tout le monde et de soulever des montagnes pour tes enfants. Merci Papa, de m'avoir transmis l'amour du détail, du travail bien fait et des 2CV. Enfin, merci Azin, d'avoir été mon ancre dans la tempête de ces dernières années.
